\section{Introduction}
\label{sec:introduction}

Wymore's Tricotyledon Theory of System Design (T3SD) provides a rigorous mathematical foundation for Model-Based Systems Engineering (MBSE) by separating what a system must do at its external boundary from how it can be constructed out of available components \citep{Wymore1993,salado2019wymore,salado2019mbse}. An Input/Output Requirement ($IOR$) specifies admissible input trajectories and their eligible output trajectories over an operational horizon, generating the \emph{functionality cotyledon} $\mathrm{CTL}(IOR)=\mathrm{FSR}(IOR)$---a space of functional system designs that is provably either empty (if the requirement violates causality) or infinite \citep{Wymore1993}. Independently, a technology ($TYR$) cataloging available hardware or software components and a System Coupling Recipe ($SCR$) wiring their ports generate the \emph{buildability cotyledon} $\mathrm{CTL}(TYR)=\mathrm{BSR}(TYR)$, whose members are coupled resultant systems $Z_b=\mathrm{RSY}(SCR)$.

Connecting these two design spaces is the central verification problem of T3SD. A candidate architecture belongs to the \emph{implementability cotyledon} $\mathrm{CTL}(IOR,TYR)=\mathrm{ISR}(IOR,TYR)$ when a buildable resultant $Z_b\in\mathrm{CTL}(TYR)$ \emph{implements} a functional system design $Z_f\in\mathrm{CTL}(IOR)$ \citep{Wymore1993}. Because a physical or digital build typically contains internal clock phases, initialization states, diagnostic registers, or unused capacity that are not part of the abstract functional model, Wymore defines implementation in two stages: one first isolates a \emph{system mode} $Z_m\in\mathrm{SYSMO}(Z_b)$ representing the operational regime of the build, and then exhibits a surjective, transition- and readout-preserving \emph{system homomorphism} $h=(h_S,h_I,h_O)$ from $Z_m$ onto $Z_f$.

While mathematically complete, this constructive workflow faces a severe practical bottleneck. First, $\mathrm{CTL}(IOR)$ is infinite when nonempty, forcing the engineer to select a particular reference machine $Z_{\mathrm{ref}}\in\mathrm{CTL}(IOR)$ and manually construct a state-space mapping $h_S:S_m\to S_{\mathrm{ref}}$ from the coupled product state space of $Z_m$ onto $Z_{\mathrm{ref}}$. Second, deciding whether a surjective homomorphism exists between finite transition systems is NP-hard in general \citep{HellNesetril1990,BodirskyKaraMartin2012}, and even small representation changes---such as flipping a phase bit or adding a diagnostic counter to a component---require reconstructing and re-proving the homomorphism maps by hand. Third, because T3SD has remained a specialized systems engineering theory, it has not developed the mature ecosystem of automated solvers and model checkers that computer science built for temporal logic \citep{Pnueli1977,BaierKatoen2008}.

This paper resolves that bottleneck by establishing a formal bridge between T3SD and temporal logic. Rather than replacing Wymore's cotyledons, we show that temporal logic provides an exact, tool-supported language for specifying the functionality cotyledon and automatically certifying homomorphic implementation of system modes. Specifically, this paper makes four contributions:
\begin{enumerate}
    \item \textbf{Isomorphism between Wymorian $IOR$s and boundary temporal specifications.} For every Wymorian $IOR$, we construct a characteristic boundary temporal specification $\Phi_{IOR}$ whose trajectory satisfaction predicate is identical to Wymore's $\mathrm{SatisfiesIOR}$, yielding a canonical bijection $\mathrm{AFSR}(\Phi_{IOR})\simeq\mathrm{FSR}(IOR)$ between the assertional functionality space and the functionality cotyledon.
    \item \textbf{Direct mode-to-$IOR$ verification without manual homomorphisms.} Using Wymore's requirement-preservation theorem (Theorem~6.58) alongside a compiled dynamics fragment $\Phi_{\mathrm{dyn}}(Z_{\mathrm{ref}})$, we prove that a system mode $Z_m\in\mathrm{SYSMO}(Z_b)$ admits an identity-boundary homomorphism onto a functional system in $\mathrm{CTL}(IOR)$ if and only if $Z_m\models\Phi_{IOR}$ directly. When a specific reference architecture $Z_{\mathrm{ref}}\in\mathrm{CTL}(IOR)$ is chosen, $Z_m\models_{\mathrm{dyn}}\Phi_{\mathrm{dyn}}(Z_{\mathrm{ref}})$ is necessary and sufficient for $Z_{\mathrm{ref}}$ to be a homomorphic image of $Z_m$, placing $(Z_{\mathrm{ref}},Z_b)$ in the implementability cotyledon $\mathrm{CTL}(IOR,TYR)$.
    \item \textbf{Decidable LTL model checking for finite buildable modes and invariance across refactoring.} Because real-world buildable components in $TYR$ are finite-state machines, any coupled resultant $Z_b=\mathrm{RSY}(SCR)$ and any operational mode $Z_m\in\mathrm{SYSMO}(Z_b)$ is a finite transition system. When $\Phi_{IOR}$ is expressed in Linear Temporal Logic (LTL), verifying $Z_m\models\Phi_{IOR}$ is decidable and fully automated via standard LTL model checking \citep{BaierKatoen2008}, and conformance is preserved under state re-encoding, port relabeling, and componentwise elaboration.
    \item \textbf{End-to-end infinite-to-finite case study on a coupled streaming Run-Length Encoder (RLE).} We develop a four-component coupled machine (tape, sense/act head, state register, and instruction table) starting from an idealized infinite-tape / unbounded-count specification in $\mathrm{CTL}(IOR_\infty)$ and systematically refining it to a finite-tape buildable resultant $Z_{b,N}=\mathrm{RSY}(SCR_N)\in\mathrm{CTL}(TYR_N)$. An automated LTL checker verifies the active-phase system mode across nominal caps, phase re-encodings, and register instrumentation while pinpointing the exact violated temporal clauses in faulty variants.
\end{enumerate}

The remainder of the paper is organized as follows. Section~\ref{sec:background} reviews the three cotyledons, system modes, homomorphisms, and temporal logic operators in accessible systems engineering terms. Section~\ref{sec:methodology} establishes the boundary $IOR$ isomorphism, the mode-to-$IOR$ implementability theorems, and the invariance and compositional lifting properties. Section~\ref{sec:case-study} presents the streaming RLE case study from its infinite-count formulation through its finite-tape refinement and automated LTL verification. Section~\ref{sec:discussion} discusses expressivity, pipeline timing, and systems engineering implications, and Section~\ref{sec:conclusions} concludes.
